% answers_short.tex —— 2004 年·普通高等学校春季招生考试·安徽卷文科 · 简短答案速查
% 由 tools/make_short_answers.py 自 answers.tex 的 \gksolhead/\gkanswer 提取，
% 并与 2004/2004_数学真题_春季安徽文_答案要点.md 逐题比对后生成，勿手改（改源 MD 后重跑）。

\begin{gknotice}
\textbf{使用说明}：本节只给最终答案便于快速核对，完整推导见解析版。答案为编写者独立推导，\textbf{非官方标准答案}。
\end{gknotice}

\gkpart{选择题}
\begin{center}\begin{tabular}{ccccccccccccc}
\toprule
题号 & 1 & 2 & 3 & 4 & 5 & 6 & 7 & 8 & 9 & 10 & 11 & 12 \\
\midrule
答案 & D & A & C & D & A & C & D & D & D & B & D & B \\
\bottomrule
\end{tabular}\end{center}

\gkpart{填空题}
\begin{itemize}[leftmargin=2.6em]
  \item[\textbf{13.}] $x=-\dfrac32$。
  \item[\textbf{14.}] $\dfrac7{10}$。
  \item[\textbf{15.}] $\dfrac14$。
  \item[\textbf{16.}] $n=3$。
\end{itemize}

\gkpart{解答题}
\begin{itemize}[leftmargin=2.6em]
  \item[\textbf{17.}] 若 $a>1$，解集为 $(1,a^2)$；若 $0<a<1$，解集为 $(a^2,1)$。
  \item[\textbf{18.}] 通项不唯一。
  \item[\textbf{19.}] (1) $k=1$ 时为圆，$k\ne1$ 时为椭圆；(2) 矩形周界，边界满足 $\max(|y|,k|x|)=\dfrac{c\sqrt{1+k^2}}2$
  \item[\textbf{20.}] (1) $\cos\angle(AC,BC_1)=\dfrac{\sqrt{10}}5$；(2) 见解析版
  \item[\textbf{21.}] 零售价定为 $30$ 元时毛利润最大，最大毛利润为 $23000$ 元。
  \item[\textbf{22.}] (1) $M=(-1,\dfrac12)$；(2) 对任意实数 $a$，顶点 $M=(-2,-\dfrac12)$ 的法线 $x=-2$ 均经过 $P$；当 $a>0$ 时，另有两点 $M=(-2\pm\sqrt a,a-\dfrac12)$，法线分别为 $y-a=\mp\dfrac{x+2}{2\sqrt a}$；当 $a\le0$ 时无其他点
\end{itemize}
